
\usepackage{listings,newtxtt}
\usepackage{xcolor}

\definecolor{commentsColor}{rgb}{0.497495, 0.497587, 0.497464}
\definecolor{keywordsColor}{rgb}{0.000000, 0.000000, 0.635294}
\definecolor{stringColor}{rgb}{0.558215, 0.000000, 0.135316}
\definecolor{arylideyellow}{rgb}{0.91, 0.84, 0.42}
\definecolor{banana}{rgb}{0.98, 0.91, 0.71}


\lstset{ %
  language=Python,                 % the language of the code (can be overrided per snippet)
  basicstyle=\ttfamily,            % fonts that are used for the code
  backgroundcolor=\color{green!10},% choose the background color; 
  breakatwhitespace=false,         % sets if automatic breaks should only happen at whitespace
  breaklines=true,                 % sets automatic line breaking
  captionpos=t,                    % sets the caption-position to bottom
  commentstyle=\color{commentsColor}\textit,    % comment style
  deletekeywords={...},            % if you want to delete keywords from the given language
  extendedchars=true,              % lets you use non-ASCII characters; does not work with UTF-8
  escapeinside={(*}{*)},           % LaTeX
  frame=t,                         % adds a frame around the code
  keepspaces=true,                 % keeps spaces in text
  keywordstyle=\color{keywordsColor}\bfseries,       % keyword style
  otherkeywords={*,...},           % if you want to add more keywords to the set
  numbers=right,                   % where to put the line-numbers; possible values are (none, left, right)
  numbersep=5pt,                   % how far the line-numbers are from the code
  %xleftmargin=.25in,               % indentação
  %xrightmargin=.25in}
  numberstyle=\tiny\color{commentsColor}, % the style that is used for the line-numbers
  rulecolor=\color{black},         % if not set, the frame-color may be changed on line-breaks within not-black text (e.g. comments (green here))
  showspaces=false,                % show spaces everywhere adding particular underscores; it overrides 'showstringspaces'
  showstringspaces=false,          % underline spaces within strings only
  showtabs=false,                  % show tabs within strings adding particular underscores
  stepnumber=1,     % the step between two line-numbers. If it's 1, each line will be numbered
  stringstyle=\color{stringColor}, % string literal style
  tabsize=1,                       % sets default tabsize to 2 spaces
  title=\lstname,                  % show the filename of files included with \lstinputlisting; also try caption instead of title
  columns=fixed                    % Using fixed column width (for e.g. nice alignment)
  %extendedchars=\true,
  %inputencoding=utf8x,
  }


\lstdefinestyle{interfaces}
  {
  float=tp,
  floatplacement=tbp,
  abovecaptionskip=-5pt
  }

\lstdefinelanguage{Python}
  {
  keywords={print, import, True, False, plot, Mesh, array, plt, def, return, type, lambda, map, list, while, if , elif, else, for, in, __name__ },
  keywordstyle=\color{blue}\bfseries,
  ndkeywords={boolean, throw, implements, import, this, not, int, float, while, open,close, read, readline, redalines,write,writelines, seek, split, len},
  ndkeywordstyle=\color{red}\bfseries,
  identifierstyle=\color{black},
  sensitive=false,
  comment=[l]{//},
  morecomment=[s]{/*}{*/},
  commentstyle=\color{purple}\ttfamily,
  stringstyle=\color{red}\ttfamily,
  morestring=[b]',
  morestring=[b]"
}



\lstset{
language=Python,   % R code
literate=	{á}{{\'a}}1
			{à}{{\`a}}1
			{ã}{{\~a}}1
{é}{{\'e}}1
{ê}{{\^e}}1
{í}{{\'i}}1
{ó}{{\'o}}1
{õ}{{\~o}}1
{ú}{{\'u}}1
{ü}{{\"u}}1
{ç}{{\c{c}}}1
}

\renewcommand{\lstlistingname}{\scriptsize Código}

\usepackage{inconsolata}
